\section{\label{sec:concs}Conclusion}

Parton distribution functions (PDFs) characterize nucleon structure in terms of 
the nucleon's constituent quarks and gluons. These PDFs are defined as matrix 
elements of light-front wave functions and cannot be directly calculated in 
Euclidean lattice QCD. In principle, however, the Mellin moments of PDFs can be 
calculated in lattice QCD, through matrix elements of twist-two operators. 
Unfortunately, these calculations are limited to the first few moments, because 
the hypercubic symmetry of the lattice regulator induces power-divergent mixing 
between twist-two operators of different mass dimension, obscuring the 
continuum limit of matrix elements determined on the lattice. Current 
determinations of PDFs rely on global analyses of data in a wide range of 
experimental channels and a determination of PDFs from first principles is 
lacking.

A new approach to determining PDFs in lattice QCD was recently proposed by Ji 
and subsequently, through a related framework, by Qiu and Ma. 
In this approach, one calculates Euclidean quasi PDFs at large nucleon 
momentum. Two recent lattice calculations 
provided promising results, but several aspects of the approach are yet to be 
fully understood. First, there is the practical issue of the systematic 
uncertainties associated with finite nucleon momenta in lattice calculations. 
This issue is likely to be resolved, to the extent that lattice 
calculations will have sufficient precision that results will provide
useful input into global analyses where experimental data are inadequate, 
with improvements in computational resources and the algorithmic advances 
already underway. Second, there are theoretical issues to be clarified: the 
renormalization of the extended operator that defines the quasi PDFs; and the 
relation between the Euclidean quasi PDF and the light-front PDF, which to date 
had been analyzed through a factorization formula at one loop in 
perturbation theory.

We have addressed the first of these theoretical considerations by introducing 
a quasi PDF 
constructed from fields smeared via the gradient flow. 
We explicitly demonstrated that there is a 
simple relation between the 
Mellin moments of the smeared Euclidean quasi PDF and the 
renormalized Mellin moments of the light-front PDF, 
once nucleon mass corrections 
are incorporated and provided the flow time is small relative to the inverse 
nucleon momentum. Corrections to this relation appear at 
${\cal O}(\Lambda_{\mathrm{QCD}}^2/P_z^2)$, where $P_z$ is the Euclidean 
momentum of the nucleon, and ${\cal O}(\sqrt{\tau} \Lambda_{\mathrm{QCD}})$, 
where 
$\tau$ is the flow time. From this correspondence it follows that, provided 
$\Lambda_{QCD},M_N\ll P_z \ll \tau^{-1/2}$, the quasi PDF and light-front PDF 
can be matched through a convolution relation. 

The chief advantage of our approach is that the gradient flow renders the quasi 
PDF finite in the continuum limit and evades the issues of the renormalization 
of 
the non-local operator that defines 
the quasi PDF on the lattice. The resulting continuum matrix elements are 
independent of the choices of discretized action used to undertake 
lattice QCD calculations and can be 
matched directly to the corresponding light-front PDFs in the $\ms$ scheme 
using 
continuum perturbation theory. Combined with a 
nonperturbative step-scaling procedure, this matching can be carried out at an 
energy sufficiently high that perturbative truncation errors are no longer 
uncontrolled. The nonperturbative implementation of our proposal is work in 
progress.
